% Generated draft: requires mathematical and source review.
\documentclass[11pt]{article}
\usepackage[margin=25mm]{geometry}
\usepackage{fontspec}
\setmainfont{Latin Modern Roman}
\usepackage{amsmath,amssymb,mathtools,unicode-math}
\setmathfont{Latin Modern Math}
\usepackage{enumitem,needspace,xurl,hyperref}
\setlength{\parindent}{0pt}
\setlength{\parskip}{5pt}
\setlength{\emergencystretch}{3em}
\newcommand{\N}{\mathbb{N}}
\newcommand{\Z}{\mathbb{Z}}
\newcommand{\Q}{\mathbb{Q}}
\newcommand{\R}{\mathbb{R}}
\newcommand{\C}{\mathbb{C}}
\newcommand{\F}{\mathbb{F}}
\newcommand{\A}{\mathbb{A}}
\newcommand{\PP}{\mathbb{P}}
\newcommand{\E}{\mathbb{E}}
\newcommand{\eps}{\varepsilon}
\newcommand{\dd}{\,\mathrm{d}}
\begin{document}
\section*{UnsolvedMath ID: 30000626}
\textbf{Source problem code:} catalogue entry
\subsection{Definitions and mathematical statement}
\textbf{Steiner-root conjecture, as supplied.} Let \(K\subset\R^n\) be a convex body, with Euclidean inradius \(r\) and circumradius \(R\). If \(a_1\le\cdots\le a_n\) are the real parts of the roots of its Steiner polynomial, counted with algebraic multiplicity, the proposed inequality is\[a_1\le-R\le-r\le a_n\le0.\] Here a convex body conventionally means compact convex with nonempty interior, and \(n\ge1\). Let \(B^n\) be the closed unit ball and define the radii by \(r=\max\{\rho:\exists c,\ c+\rho B^n\subseteq K\}\) and \(R=\min\{\rho:\exists c,\ K\subseteq c+\rho B^n\}\). The centers may differ; no concentric-ball assumption is imposed.\par Under the standard Euclidean normalization the Steiner polynomial \(S_K(z)\) is the degree-\(n\) polynomial determined by\[S_K(t)=\operatorname{vol}_n(K+tB^n)\quad(t\ge0).\] Its complex roots are \(\zeta_1,\ldots,\zeta_n\), and the \(a_i\) are the sorted multiset of \(\operatorname{Re}\zeta_i\). The conjecture does not assert that all roots are real or order the complex roots themselves. The chain says that the smallest real part reaches at least as far left as \(-R\), the largest reaches at least as far right as \(-r\), and no root has positive real part. It does not place every root real part inside \([-R,-r]\); reversing those endpoint bounds would change the claim. The middle comparison follows the radius order and is retained exactly.\par The source supplies no dimension restriction or additional symmetry, smoothness, or reference-body hypothesis, so none is added. The supplied triage checked 2026-08-21 labels the chain open and reports no located complete resolution, with the report as its sole source. That dated assertion is not independently applied here as a verified current status in every dimension.
\subsection{Short English statement}
The proposed chain bounds the extreme real parts of the Euclidean Steiner roots using the two radii and also forbids positive real parts. It is not a claim that all roots are real or lie between the two negative radii. The full all-dimensions wording is preserved with an unverified dated status label.
\subsection{Sources}
\begin{itemize}\item Local UnsolvedMath registry, permanent ID 30000626, code catalogue entry; \emph{Root Bounds for Steiner Polynomials of Convex Bodies}, Oberwolfach Reports Open Problems.\item \emph{Konvexgeometrie (2007)}, Oberwolfach Reports, DOI 10.4171/owr/2006/56, pp. 3321--3396. \url{https://doi.org/10.4171/owr/2006/56}. Proposed year and year tag: 2006; no proposer named.\item The same full report is the sole supplied literature source. No separate authoritative root-bound theorem or later status reference is supplied.\item Literature triage checked 2026-08-21, status \texttt{open}; record updated 2026-09-16. The \texttt{corrected\_verified} field records notation normalization and preservation of the full chain. Verify the dated unresolved-status claim independently.\end{itemize}
\end{document}
